\documentclass{article}
\usepackage{amsmath}
\usepackage{amssymb}

\title{Advanced LaTeX Math Environments}
\author{htex Library}
\date{2026}

\begin{document}

\maketitle

\section{Piecewise Functions with Cases}

The absolute value function can be defined using the \texttt{cases} environment:

\begin{equation}
|x| = \begin{cases}
  -x & \text{if } x < 0 \\
  x & \text{if } x \geq 0
\end{cases}
\end{equation}

The Heaviside step function is another common use:

\begin{equation}
H(t) = \begin{cases}
  0 & t < 0 \\
  1/2 & t = 0 \\
  1 & t > 0
\end{cases}
\end{equation}

\section{Multi-line Equations with Split}

When you need to align multi-line equations without explicit alignment markers:

\begin{equation}
\begin{split}
(x + y)^3 &= x^3 + 3x^2y + 3xy^2 + y^3 \\
&= x^3 + 3xy(x + y) + y^3
\end{split}
\end{equation}

\section{Aligned Equations}

Use the \texttt{aligned} environment for inline alignment:

\begin{equation}
\begin{aligned}
2x + 3y &= 7 \\
x - y &= 1
\end{aligned}
\end{equation}

\section{Matrix Varieties}

\subsection{Standard Matrix (no delimiters)}

\begin{equation}
A = \begin{matrix}
  1 & 2 & 3 \\
  4 & 5 & 6 \\
  7 & 8 & 9
\end{matrix}
\end{equation}

\subsection{Parenthesis Matrix}

\begin{equation}
B = \begin{pmatrix}
  a & b \\
  c & d
\end{pmatrix}
\end{equation}

\subsection{Square Bracket Matrix}

\begin{equation}
C = \begin{bmatrix}
  1 & 0 & 0 \\
  0 & 1 & 0 \\
  0 & 0 & 1
\end{bmatrix}
\end{equation}

\subsection{Determinant (Vertical Bar Matrix)}

\begin{equation}
\det(A) = \begin{vmatrix}
  a & b \\
  c & d
\end{vmatrix} = ad - bc
\end{equation}

\subsection{Double Vertical Bar Matrix (Norm)}

\begin{equation}
\|v\| = \begin{Vmatrix}
  v_1 \\
  v_2 \\
  v_3
\end{Vmatrix}
\end{equation}

\subsection{Small Matrix (inline)}

The matrix $\begin{smallmatrix} 1 & 2 \\ 3 & 4 \end{smallmatrix}$ can fit inline.

\section{Gathered Equations}

For centered, multi-line equations without alignment:

\begin{equation}
\begin{gathered}
f(x) = x^2 + 2x + 1 \\
= (x + 1)^2
\end{gathered}
\end{equation}

\section{Aligned with Fixed Columns}

Using \texttt{alignedat} for multiple alignment points:

\begin{equation}
\begin{alignedat}{2}
x &= a + b \quad & y &= c + d \\
x - y &= a - c \quad & 2x &= a + b + c + d
\end{alignedat}
\end{equation}

\section{Subarray for Indices}

Use \texttt{subarray} for multi-line superscripts or subscripts:

\begin{equation}
\sum_{\begin{subarray}{l}
  i=1 \\
  j=1
\end{subarray}}^{n} a_{ij}
\end{equation}

\section{Display Cases}

The \texttt{dcases} environment uses display-size fonts:

\begin{equation}
f(x) = \begin{dcases}
  \frac{1}{2}(x^2 + 1) & \text{if } |x| \leq 1 \\
  \frac{1}{2}(|x| + 1) & \text{if } |x| > 1
\end{dcases}
\end{equation}

\end{document}
